\ifdefined\gkver\else\def\gkver{student}\fi
\documentclass[\gkver]{gaokaozhenti}
\begin{document}

\chapter{2003年上海}

\begin{enumerate}
\item
%% number: 1
%% paperName: 2003年上海高考第 1 题 5 分
%% typeId: 1
%% type: 单选题
%% chapter: 原子和原子核
%% point: 人工核反应
%% method: 无
%% score: 5
%% degree: 850
%% duplicateId: 0
%% body:
在核反应方程 $\ce{^{4}_{2}He}+\ce{^{14}_{7}N}\to\ce{^{17}_{8}O}+$（X）的括弧中，X 所代表的粒子是\xzanswer{A}

\fourchoices[answer=A]
{$\ce{^{1}_{1}H}$}
{$\ce{^{2}_{1}H}$}
{$\ce{^{0}_{-1}e}$}
{$\ce{^{1}_{0}n}$}

%% answer: A
%% memo:
\memoanswer{
核反应方程满足质量数和核电荷数守恒，由核反应方程知，粒子的质量数为 1，核电荷数为 1，故 A 正确。
}

\item
%% number: 2
%% paperName: 2003年上海高考第 2 题 5 分
%% typeId: 2
%% type: 多选题
%% chapter: 机械波
%% point: 机械波的产生
%% method: 无
%% score: 5
%% degree: 850
%% duplicateId: 0
%% body:
关于机械波，下列说法正确的是\xzanswer{ABC}

\fourchoices[answer=ABC]
{在传播过程中能传递能量}
{频率由波源决定}
{能产生干涉、衍射现象}
{能在真空中传播}

%% answer: ABC
%% memo:
\memoanswer{
机械波的频率由波源的振动频率决定，传播的过程中需要介质，且能产生干涉和衍射现象，故 ABC 正确，D 错误。
}

\item
%% number: 3
%% paperName: 2003年上海高考第 3 题 5 分
%% typeId: 1
%% type: 单选题
%% chapter: 光学
%% point: 光电效应
%% method: 无
%% score: 5
%% degree: 850
%% duplicateId: 0
%% body:
爱因斯坦由光电效应的实验规律，猜测光具有粒子性，从而提出光子说，从科学的方法来说，这属于\xzanswer{C}

\fourchoices[answer=C]
{等效替代}
{控制变量}
{科学假说}
{数学归纳}

%% answer: C
%% memo:
\memoanswer{
根据实验现象，而后提出理论支持，物理学上把这种研究方法称为“科学假说”，故 C 正确。
}

\item
%% number: 4
%% paperName: 2003年上海高考第 4 题 5 分
%% typeId: 2
%% type: 多选题
%% chapter: 功和能
%% point: 动能定理
%% method: 无
%% score: 5
%% degree: 650
%% duplicateId: 0
%% body:
一个质量为 $0.3\Ukg$ 的弹性小球，在光滑水平面上以 $6\Ums$ 的速度垂直撞到墙上，碰撞后小球沿相反方向运动，反弹后的速度大小与碰撞前相同，则碰撞前后小球速度变化量的大小 $\Delta v$ 和碰撞过程中墙对小球做功的大小 $W$ 为\xzanswer{BC}

\fourchoices[answer=BC]
{$\Delta v=0$}
{$\Delta v=12\Ums$}
{$W=0$}
{$W=10.8\UJ$}

%% answer: BC
%% memo:
\memoanswer{
速度是矢量，速度的变化也是矢量，则速度变化量 $\Delta v=-v-v=-2v$，其大小为 $2v=12\Ums$，反弹前后小球的动能没有变化，$\Delta E_{k}=0$，则墙对小球做功的大小 $W=\Delta E_{k}=0$，故 BC 正确。
}

\item
%% number: 5
%% paperName: 2003年上海高考第 5 题 5 分
%% typeId: 2
%% type: 多选题
%% chapter: 电场
%% point: 电势能电势差电场力做功
%% method: 无
%% score: 5
%% degree: 650
%% duplicateId: 0
%% body:
一负电荷仅受电场力作用，从电场中的 A 点运动到 B 点，在此过程中该电荷作初速度为零的匀加速直线运动，则 A、B 两点电场强度 $E_{A}$、$E_{B}$ 及该电荷在 A、B 两点的电势能 $E_{pA}$、$E_{pB}$ 之间的关系为\xzanswer{AD}

\fourchoices[answer=AD]
{$E_{A}=E_{B}$}
{$E_{A}<E_{B}$}
{$E_{pA}=E_{pB}$}
{$E_{pA}>E_{pB}$}

%% answer: AD
%% memo:
\memoanswer{
因负电荷做匀加速直线运动，可知电场为匀强电场，负电荷逐渐加速，则电场力做正功，动能增大，电势能减小，故 AD 正确。
}

\item
%% number: 6
%% paperName: 2003年上海高考第 6 题 5 分
%% typeId: 1
%% type: 单选题
%% chapter: 电磁感应
%% point: 电磁感应电路分析
%% method: 无
%% score: 5
%% degree: 450
%% duplicateId: 0
%% body:
粗细均匀的电阻丝围成的正方形线框置于有界匀强磁场中，磁场方向垂直于线框平面，其边界与正方形线框的边平行，现使线框以同样大小的速度沿四个不同方向平移出磁场，如图所示，则在移动过程中线框的一边 a、b 两点间电势差绝对值最大的是\xzanswer{B}

\fourchoices[answer=B, ispicture=true, wA=2.3cm, wB=2.3cm, wC=3.2cm, wD=3.2cm]
{figs/06a.png}
{figs/06b.png}
{figs/06c.png}
{figs/06d.png}

%% answer: B
%% memo:
\memoanswer{
在四个图中，线框切割产生的电动势大小均相等（$E$），回路电阻均为 $4r$，则电路中电流亦相等（$I$），B 图中，ab 为电源，$U_{ab}$ 为路端电压，故 $U_{ab}=I\cdot3r=\frac{3}{4}E$，其他情况下，$U_{ab}=I\cdot r=\frac{1}{4}E$，故 B 正确。
}

\item
%% number: 7
%% paperName: 2003年上海高考第 7 题 5 分
%% typeId: 2
%% type: 多选题
%% chapter: 功和能
%% point: 整体机械能守恒
%% method: 极值
%% score: 5
%% degree: 450
%% duplicateId: 0
%% body:
一质量不计的直角形支架两端分别连接质量为 $m$ 和 $2m$ 的小球 A 和 B，支架的两直角边长度分别为 $2L$ 和 $L$，支架可绕固定轴 O 在竖直平面内无摩擦转动，如图所示，开始时 OA 边处于水平位置，由静止释放，则\xzanswer{BCD}

\onepicture[w=4.6cm]{\includesvg{figs/07}}

\fourchoices[answer=BCD]
{A 球的最大速度为 $2\sqrt{gL}$}
{A 球速度最大时，两小球的总重力势能最小}
{A 球速度最大时，两直角边与竖直方向的夹角为 $45^{\circ}$}
{A、B 两球的最大速度之比 $v_{A}:v_{B}=2:1$}

%% answer: BCD
%% memo:
\memoanswer{
设 A 球的速度为 $v_{A}$，B 球的速度为 $v_{B}$，A、B 运动时角速度相等，则 $v_{A}:v_{B}=\omega r_{A}:\omega r_{B}=r_{A}:r_{B}=2:1$，系统动能最大时，OB 与竖直方向的夹角为 $\alpha$，由机械能守恒得，则有 $mg\cdot2L\sin\alpha-2mg\cdot L(1-\cos\alpha)=\frac{1}{2}mv_{A}^{2}+\frac{1}{2}\cdot2mv_{B}^{2}$，A 球的动能最大时，系统的动能最大，系统的势能最小，化简得 $v_{A}^{2}=\frac{8gL}{3}[\sqrt{2}\sin(\alpha+45^{\circ})-1]$，得 $\alpha=45^{\circ}$ 时取 $v_{A}^{2}$ 最大值，则 $v_{A\max}^{2}=\frac{8gL}{3}(\sqrt{2}-1)$，即 $v_{A\max}=\sqrt{\frac{8gL}{3}(\sqrt{2}-1)}$，故 A 错误 BCD 正确。
}

\item
%% number: 8
%% paperName: 2003年上海高考第 8 题 5 分
%% typeId: 1
%% type: 单选题
%% chapter: 光学
%% point: 光的干涉
%% method: 无
%% score: 5
%% degree: 450
%% duplicateId: 0
%% body:
劈尖干涉是一种薄膜干涉，其装置如图\ref{2003上海08a}所示，将一块平板玻璃放置在另一平板玻璃之上，在一端夹入两张纸片，从而在两片玻璃表面之间形成一个劈形空气薄膜，当光垂直入射后，从上往下看到的干涉条纹如图\ref{2003上海08b}所示，干涉条纹有如下特点：

（1）任意一条明条纹或暗条纹所在位置下面的薄膜厚度相等；

（2）任意相邻明条纹或暗条纹所对应的薄膜厚度差恒定。

现若将图\ref{2003上海08a}装置中抽去一张纸片，则当光垂直入射到新的劈形空气膜后，从上往下观察到的干涉条纹\xzanswer{A}

\twopicture[numstyle=arabic, w=5.5cm, capA={俯视图}, labelA=2003上海08a, labelB=2003上海08b]
{figs/08a.png}
{figs/08b.png}

\fourchoices[answer=A]
{变疏}
{变密}
{不变}
{消失}

%% answer: A
%% memo:
\memoanswer{
撤去一张纸片，相同水平距离上，劈尖厚度减小，干涉时波程差减小，所以条纹变宽，数量变少，故 A 正确。
}

\item
%% number: 9
%% paperName: 2003年上海高考第 9 题 4 分
%% typeId: 3
%% type: 填空题
%% chapter: 原子和原子核
%% point: 原子核式结构
%% method: 无
%% score: 4
%% degree: 650
%% duplicateId: 0
%% body:
卢瑟福通过\tkanswer[2.6cm]{$\alpha$ 粒子散射}实验，发现了原子中间有一个很小的核，并由此提出了原子的核式结构模型，右面平面示意图中的四条线表示 $\alpha$ 粒子运动的可能轨迹，在图中完成中间两条 $\alpha$ 粒子的运动轨迹。

\onepicture[align=right, w=5.2cm]{\includesvg{figs/09}}

%% answer:
%% memo:
\memoanswer{
$\alpha$ 粒子散射；运动轨迹如图所示。

\onepicture[w=4.8cm]{\includesvg{figs/09_an}}
}

\item
%% number: 10
%% paperName: 2003年上海高考第 10 题 4 分
%% typeId: 3
%% type: 填空题
%% chapter: 机械波
%% point: 波的图象
%% method: 无
%% score: 4
%% degree: 650
%% duplicateId: 0
%% body:
细绳的一端在外力作用下从 $t=0$ 时刻开始做简谐振动，激发出一列简谐波，在细绳上选取 15 个点，图 1 为 $t=0$ 时刻各质点所处的位置，图 2 为 $t=\frac{T}{4}$ 时刻的波形图（$T$ 为波的周期），在图 3 中画出 $t=\frac{3T}{4}$ 时刻的波形图。

\onepicture[w=13.0cm]{\includesvg{figs/10}}

%% answer:
\jdanswer{波形图如图所示。}
%% memo:
\memoanswer{
波形图如图所示。

\onepicture[w=12.0cm]{\includesvg{figs/10_an}}
}

\item
%% number: 11
%% paperName: 2003年上海高考第 11 题 4 分
%% typeId: 3
%% type: 填空题
%% chapter: 曲线运动
%% point: 万有引力定律
%% method: 类比
%% score: 4
%% degree: 650
%% duplicateId: 0
%% body:
有质量的物体周围存在着引力场，万有引力和库仑力有类似的规律，因此我们可以用定义静电场场强的方法来定义引力场的场强，由此可得，与质量为 $M$ 的质点相距 $r$ 处的引力场场强的表达式为 $E_{G}=$\tkanswer[2.2cm]{$\frac{GM}{r^{2}}$}（万有引力恒量用 $G$ 表示）。

%% answer:
%% memo:
\memoanswer{
$\frac{GM}{r^{2}}$。
}

\item
%% number: 12
%% paperName: 2003年上海高考第 12 题 4 分
%% typeId: 1
%% type: 单选题
%% chapter: 电场
%% point: 带电粒子的圆周运动
%% method: 等效替代
%% score: 4
%% degree: 650
%% duplicateId: 0
%% body:
若氢原子的核外电子绕核作半径为 $r$ 的匀速圆周运动，则其角速度 $\omega=$\tkanswer[2.2cm]{$\frac{e}{r}\sqrt{\frac{k}{mr}}$}，电子绕核运动可等效为环形电流，则电子运动的等效电流 $I=$\tkanswer[2.8cm]{$\frac{e^{2}}{2\pi r}\sqrt{\frac{k}{mr}}$}（已知电子的质量为 $m$，电量为 $e$，静电力恒量用 $k$ 表示）。

%% answer:
%% memo:
\memoanswer{
由题意知电子轨道半径为 $r$，角速度为 $\omega$，原子核对电子的库仑力提供向心力，由牛顿第二定律和库仑定律得 $\frac{ke^{2}}{r^{2}}=m\omega^{2}r$，解得 $\omega=\frac{e}{r}\sqrt{\frac{k}{mr}}$，又 $T=\frac{2\pi}{\omega}$，由电流的定义得 $I=\frac{q}{t}=\frac{e}{T}$，联立解得 $I=\frac{e^{2}}{2\pi r}\sqrt{\frac{k}{mr}}$。
}

\item
%% number: 13
%% paperName: 2003年上海高考第 13 题 4 分
%% typeId: 3
%% type: 填空题
%% chapter: 气体性质
%% point: 查理定律
%% method: 无
%% score: 4
%% degree: 450
%% duplicateId: 0
%% body:
某登山爱好者在攀登珠穆朗玛峰的过程中，发现他携带的手表表面玻璃发生了爆裂，这种手表是密封的，出厂时给出的参数为：$27\UduC$ 时表内气体压强为 $1\times10^{5}\UPa$，在内外压强差超过 $6\times10^{4}\UPa$ 时，手表表面玻璃可能爆裂，已知当时手表处的气温为 $-13\UduC$，则手表表面玻璃爆裂时表内气体压强的大小为\tkanswer[2.4cm]{$8.7\times10^{4}$}$\UPa$，已知外界大气压强随高度变化而变化，高度每上升 $12\Um$ 大气压强降低 $133\UPa$，设海平面大气压为 $1\times10^{5}\UPa$，则登山运动员此时的海拔高度约为\tkanswer[1.8cm]{$6613$}$\Um$。

%% answer:
%% memo:
\memoanswer{
手表内封闭气体的初始状态参量为 $p_{0}=1\times10^{5}\UPa$，$T_{0}=(273+27)\UK=300\UK$，末状态参量为 $T_{1}=(273-13)\UK=260\UK$，由查理定律得 $\frac{p_{0}}{T_{0}}=\frac{p_{1}}{T_{1}}$，解得 $p_{1}=\frac{T_{1}}{T_{0}}p_{0}=\frac{260}{300}\times10^{5}\UPa\approx8.7\times10^{4}\UPa$；手表表面玻璃爆裂时表内气体压强为 $p_{1}=8.7\times10^{4}\UPa$，由题意知，此时表外大气压强为 $p'=p_{1}-6\times10^{4}\UPa=2.7\times10^{4}\UPa$，因海平面大气压为 $1\times10^{5}\UPa$，则因为海拔升高引起大气压强的变化量为 $\Delta p=1\times10^{5}\UPa-2.7\times10^{4}\UPa=7.3\times10^{4}\UPa$，则海拔高度 $H=\frac{\Delta p}{133\UPa}\times12\Um=\frac{7.3\times10^{4}}{133}\times12\Um=6586\Um$。
}

\item
%% number: 14
%% paperName: 2003年上海高考第 14 题 5 分
%% typeId: 4
%% type: 实验题
%% chapter: 曲线运动
%% point: 研究平抛运动实验
%% method: 无
%% score: 5
%% degree: 650
%% duplicateId: 0
%% body:
如图所示，在研究平抛运动时，小球 A 沿轨道滑下，离开轨道末端（末端水平）时撞开轻质接触式开关 S，被电磁铁吸住的小球 B 同时自由下落，改变整个装置的高度 $H$ 做同样的实验，发现位于同一高度的 A、B 两球总是同时落地，该实验现象说明了 A 球在离开轨道后\xzanswer{C}

\onepicture[align=right, w=4.6cm]{\includesvg{figs/14}}

\fourchoices[answer=C]
{水平方向的分运动是匀速直线运动}
{水平方向的分运动是匀加速直线运动}
{竖直方向的分运动是自由落体运动}
{竖直方向的分运动是匀速直线运动}

%% answer: C
\jdanswer{C}
%% memo:
\memoanswer{
因为 A、B 两球同时从同一高度开始下落，并且同时到达地面，故在竖直方向两球遵循相同的运动规律，即速度、加速度总是相同。因 B 球做自由落体运动，故 A 球在平抛过程中在竖直方向也做自由落体运动，故 C 正确 D 错误；而 A 球在水平方向的运动没有可以参照的物体，故无法确定平抛运动的物体在水平方向的运动所遵循的规律，故 AB 错误。
}

\item
%% number: 15
%% paperName: 2003年上海高考第 15 题 5 分
%% typeId: 4
%% type: 实验题
%% chapter: 光学
%% point: 光电效应
%% method: 无
%% score: 5
%% degree: 650
%% duplicateId: 0
%% body:
在右图所示的光电管的实验中，发现用一定频率的 A 单色光照射光电管时，电流表指针会发生偏转，而用另一频率的 B 单色光照射时不发生光电效应，那么\xzanswer{AC}

\onepicture[align=right, w=4.4cm]{\includesvg{figs/15}}

\fourchoices[answer=AC]
{A 光的频率大于 B 光的频率}
{B 光的频率大于 A 光的频率}
{用 A 光照射光电管时流过电流表 G 的电流方向是 a 流向 b}
{用 A 光照射光电管时流过电流表 G 的电流方向是 b 流向 a}

%% answer: AC
\jdanswer{AC}
%% memo:
\memoanswer{
用一定频率的 A 单色光照射光电管时，电流表指针会发生偏转，知 $v_{A}>v_{0}$；用另一频率的 B 单色光照射时不发生光电效应，知 $v_{B}<v_{0}$，所以 A 光的频率大于 B 的频率，故 A 正确 B 错误；发生光电效应时，电子从光电管右端运动到左端，而电流的方向与电子定向移动的方向相反，所以流过电流表 G 的电流方向是 a 流向 b，故 C 正确 D 错误。
}

\item
%% number: 16
%% paperName: 2003年上海高考第 16 题 6 分
%% typeId: 4
%% type: 实验题
%% chapter: 力物体的平衡
%% point: 力矩盘
%% method: 无
%% score: 6
%% degree: 650
%% duplicateId: 0
%% body:
如图所示，在“有固定转动轴物体的平衡条件”实验中，调节力矩盘使其平衡，弹簧秤的读数为\tkanswer[1.4cm]{$1.9$}$\UN$，此时力矩盘除受到钩码作用力 $F_{1}$、$F_{2}$、$F_{3}$ 和弹簧秤拉力 $F_{4}$ 外，主要还受到\tkanswer[0.9cm]{重}力和\tkanswer[1.2cm]{支持}力的作用；如果每个钩码的质量均为 $0.1\Ukg$，盘上各圆的半径分别是 $0.05\Um$、$0.10\Um$、$0.15\Um$、$0.20\Um$（取 $g=10\Umsq$），则 $F_{2}$ 的力矩是\tkanswer[1.2cm]{$0.1$}$\UNcdotm$。有同学在做这个实验时，发现顺时针力矩之和与逆时针力矩之和存在较大差距，检查发现读数和计算均无差错，请指出造成这种差距的一个可能原因，并提出简单的检验方法（如例所示，将答案填在下表空格中）。

\begin{center}
\begin{tblr}{colspec={|c|X[c]|X[l]|}, width=\linewidth, hlines}
 & 可能原因 & 检验方法 \\
例 & 力矩盘面没有调到竖直 & 用一根细线挂一个钩码靠近力矩盘面，如果细线与力矩盘面间存在一个小的夹角，说明力矩盘不竖直 \\
答 &  &  \\
\end{tblr}
\end{center}

\onepicture[align=right, w=6.0cm]{\includesvg{figs/16}}

%% answer:
%% memo:
\memoanswer{
弹簧秤的读数为 $1.9\UN$（1.8～2.0 均可）；重；支持；$F_{2}$ 的力矩是 $0.1\UNcdotm$。转轴摩擦力太大，安装力矩盘后，轻轻转动盘面，如果盘面转动很快停止，说明摩擦太大；力矩盘重心没在中心，在最低端放一标志，轻轻转动盘面，如果标志始终停留在最低端，说明重心在这个标志和中心之间。
}

\item
%% number: 17
%% paperName: 2003年上海高考第 17 题 7 分
%% typeId: 4
%% type: 实验题
%% chapter: 气体性质
%% point: 验证玻意耳定律
%% method: 无
%% score: 7
%% degree: 650
%% duplicateId: 0
%% body:
有同学在做“研究温度不变时气体的压强跟体积的关系”实验时，用连接计算机的压强传感器直接测得注射器内气体的压强值，缓慢推动活塞，使注射器内空气柱从初始体积 $20.0\UmL$ 变为 $12.0\UmL$，实验共测了五次，每次体积值直接从注射器的刻度上读出并输入计算机，同时由压强传感器测得对应体积的压强值，实验完成后，计算机屏幕上立刻显示出如下表中所示的实验结果。

\begin{table}[htbp]
\centering
\begin{tblr}{|c|c|c|c|}
\hline
序号 & $V$（$\UmL$） & $p$（$\times10^{5}\UPa$） & $pV$（$\times10^{5}\UPa\cdot\UmL$） \\
\hline
1 & 20.0 & 1.0010 & 20.020 \\
\hline
2 & 18.0 & 1.0952 & 19.714 \\
\hline
3 & 16.0 & 1.2313 & 19.701 \\
\hline
4 & 14.0 & 1.4030 & 19.642 \\
\hline
5 & 12.0 & 1.6351 & 19.621 \\
\hline
\end{tblr}
\end{table}

\begin{enumerate}
	\item
	仔细观察不难发现，$pV$（$\times10^{5}\UPa\cdot\UmL$）一栏中的数值越来越小，造成这一现象的可能原因是\xzanswer{D}

	\fourchoices[answer=D]
	{实验时注射器活塞与筒壁间的摩擦力不断增大}
	{实验时环境温度增大了}
	{实验时外界大气压强发生了变化}
	{实验时注射器内的空气向外发生了泄漏}
	\item
	根据你在（1）中的选择，说明为了减小误差，应采取的措施是：\tkanswer[5.2cm]{在注射器活塞上涂上润滑油增加密封性}。
\end{enumerate}

\onepicture[align=right, w=5.6cm]{\includesvg{figs/17}}

%% answer:
\jdanswer{
\begin{enumerate}
	\item
	D
	\item
	在注射器活塞上涂上润滑油增加密封性
\end{enumerate}
}
%% memo:
\memoanswer{
（1）D。

（2）在注射器活塞上涂上润滑油增加密封性。
}

\item
%% number: 18
%% paperName: 2003年上海高考第 18 题 7 分
%% typeId: 4
%% type: 实验题
%% chapter: 电路
%% point: 分压与分流
%% method: 无
%% score: 7
%% degree: 650
%% duplicateId: 0
%% body:
图\ref{2003上海18a}为某一热敏电阻（电阻值随温度的改变而改变，且对温度很敏感）的 $I-U$ 关系曲线图。

\begin{enumerate}
	\item
	为了通过测量得到图\ref{2003上海18a}所示 $I-U$ 关系的完整曲线，在图\ref{2003上海18b}和图\ref{2003上海18c}两个电路中应选择的是图\tkanswer[0.8cm]{2}，简要说明理由：\tkanswer{图\ref{2003上海18b}电路电压可从 $0\UV$ 调到所需电压，调节范围较大（或图\ref{2003上海18c}电路不能测得 $0\UV$ 附近的数据）}。（电源电动势为 $9\UV$，内阻不计，滑线变阻器的阻值为 $0\sim100\UO$）
	\item
	在图\ref{2003上海18d}电路中，电源电压恒为 $9\UV$，电流表读数为 $70\UmA$，定值电阻 $R_{1}=250\UO$，由热敏电阻的 $I-U$ 关系曲线可知，热敏电阻两端的电压为\tkanswer[1.2cm]{$5.2$}$\UV$，电阻 $R_{2}$ 的阻值为\tkanswer[1.6cm]{$111.8$}$\UO$。
	\item
	举出一个可以应用热敏电阻的例子：\tkanswer[3.0cm]{热敏温度计}。
\end{enumerate}

\fourpicture[numstyle=arabic, h=3.4cm, align=right, labelA=2003上海18a, labelB=2003上海18b, labelC=2003上海18c, labelD=2003上海18d]
{\includesvg{figs/18a}}
{\includesvg{figs/18b}}
{\includesvg{figs/18c}}
{\includesvg{figs/18d}}

%% answer:
\jdanswer{
\begin{enumerate}
	\item
	2，图\ref{2003上海18b}电路电压可从 $0\UV$ 调到所需电压，调节范围较大（或图\ref{2003上海18c}电路不能测得 $0\UV$ 附近的数据）
	\item
	$5.2\UV$，$111.8\UO$（111.6～112.0 均给分）
	\item
	热敏温度计（提出其它实例，只要合理的都给分）
\end{enumerate}
}
%% memo:
\memoanswer{
（1）图中的 $I-U$ 关系图线是 $0\sim6\UV$ 间连续曲线，且热敏电阻的阻值非定值，所以在选择实验电路时电路必须能保证提供 $0\sim6\UV$ 间的连续电压，图 2 为分压电路，电压可从 $0\UV$ 开始调节，而图 3 电路不能测得 $0\UV$ 附近的数据。因此选择图 2。

（2）定值电阻的阻值 $R_{1}=250\UO$，两端电压为 $9\UV$，由欧姆定律可以算出流过定值电阻 $R_{1}$ 的电流为 $36\UmA$，故流过热敏电阻的电流为 $34\UmA$，由图线查出这时热敏电阻两端的电压为 $5.2\UV$，那么电阻 $R_{2}$ 两端的电压为 $3.8\UV$，结合电流可以算出它的电阻约为 $111.8\UO$。

（3）可用于制作热敏温度计及温控敏感元件等。
}

\item
%% number: 19
%% paperName: 2003年上海高考第 19 题 10 分
%% typeId: 6
%% type: 计算题
%% chapter: 气体性质
%% point: 理想气体状态方程
%% method: 无
%% score: 10
%% degree: 650
%% duplicateId: 0
%% body:
如图所示，1、2、3 为 $p-V$ 图中一定量理想气体的三个状态，该理想气体由状态 1 经过程 $1\to3\to2$ 到达状态 2，试利用气体实验定律证明：$\frac{p_{1}V_{1}}{T_{1}}=\frac{p_{2}V_{2}}{T_{2}}$。

\onepicture[align=right, w=4.6cm]{\includesvg{figs/19}}

%% answer:
\jdanswer{$\frac{p_{1}V_{1}}{T_{1}}=\frac{p_{2}V_{2}}{T_{2}}$}
%% memo:
\memoanswer{
设状态 3 的温度为 $T$。$1\to3$ 为等压过程，有 $\frac{V_{1}}{T_{1}}=\frac{V_{2}}{T}$；$3\to2$ 为等容过程，有 $\frac{p_{1}}{T}=\frac{p_{2}}{T_{2}}$。消去 $T$ 即得 $\frac{p_{1}V_{1}}{T_{1}}=\frac{p_{2}V_{2}}{T_{2}}$。
}

\item
%% number: 20
%% paperName: 2003年上海高考第 20 题 10 分
%% typeId: 5
%% type: 辨析题
%% chapter: 曲线运动
%% point: 平抛运动
%% method: 无
%% score: 10
%% degree: 650
%% duplicateId: 0
%% body:
如图所示，一高度为 $h=0.2\Um$ 的水平面在 A 点处与一倾角为 $\theta=30^{\circ}$ 的斜面连接，一小球以 $v_{0}=5\Ums$ 的速度在平面上向右运动。求小球从 A 点运动到地面所需的时间（平面与斜面均光滑，取 $g=10\Umsq$）。某同学对此题的解法为：

小球沿斜面运动，则 $\frac{h}{\sin\theta}=v_{0}t+\frac{1}{2}g\sin\theta\cdot t^{2}$，由此可求得落地时间 $t$。

问：你同意上述解法吗？若同意，求出所需时间；

若不同意则说明理由并求出你认为正确的结果。

\onepicture[align=right, w=5.6cm]{\includesvg{figs/20}}

%% answer:
\jdanswer{不同意。$t=0.2\Us$。}
%% memo:
\memoanswer{
不同意，小球应在 A 点离开平面做平抛运动，而不是沿斜面下滑，正确做法是：

假设落地点在平面上，由平抛规律得 $h=\frac{1}{2}gt^{2}$，解得 $t=\sqrt{\frac{2h}{g}}$。

落地点与 A 点的水平距离为 $s=v_{0}t=v_{0}\sqrt{\frac{2h}{g}}=1\Um$。

斜面底宽 $l=h\cot\theta=0.2\times\sqrt{3}\Um=0.35\Um$。

又 $l<s$，故小球离开 A 点后不会落到斜面上，因此落地时间即为平抛运动时间，有 $t=\sqrt{\frac{2h}{g}}=\sqrt{\frac{2\times0.2}{10}}\Us=0.2\Us$。
}

\item
%% number: 21
%% paperName: 2003年上海高考第 21 题 12 分
%% typeId: 6
%% type: 计算题
%% chapter: 功和能
%% point: 动能定理
%% method: 无
%% score: 12
%% degree: 650
%% duplicateId: 0
%% body:
质量为 $m$ 的飞机以水平速度 $v_{0}$ 飞离跑道后逐渐上升，若飞机在此过程中水平速度保持不变，同时受到重力和竖直向上的恒定升力（该升力由其它力的合力提供，不含重力），今测得当飞机在水平方向的位移为 $l$ 时，它的上升高度为 $h$，求：

\begin{enumerate}
	\item
	飞机受到的升力的大小；
	\item
	从起飞到上升到 $h$ 高度的过程中升力所作的功及高度 $h$ 处飞机的动能。
\end{enumerate}

\onepicture[align=right, w=5.0cm]{\includesvg{figs/21}}

%% answer:
\jdanswer{
\begin{enumerate}
	\item
	$F=mg\left(1+\frac{2hv_{0}^{2}}{gl^{2}}\right)$
	\item
	$W=mgh\left(1+\frac{2hv_{0}^{2}}{gl^{2}}\right)$，$E_{k}=\frac{1}{2}mv_{0}^{2}\left(1+\frac{4h^{2}}{l^{2}}\right)$
\end{enumerate}
}
%% memo:
\memoanswer{
（1）飞机水平方向速度不变，有 $l=v_{0}t$。

竖直方向加速度恒定，有 $h=\frac{1}{2}at^{2}$，消去 $t$ 得 $a=\frac{2hv_{0}^{2}}{gl^{2}}$。

由牛顿第二定律得 $F=mg+ma=mg\left(1+\frac{2hv_{0}^{2}}{gl^{2}}\right)$。

（2）升力做功为 $W=Fh=mgh\left(1+\frac{2hv_{0}^{2}}{gl^{2}}\right)$。

在 $h$ 处，飞机竖直方向的速度为 $v_{t}=\sqrt{2ah}=\frac{2hv_{0}}{l}$。

所以 $E_{k}=\frac{1}{2}m(v_{0}^{2}+v_{t}^{2})=\frac{1}{2}mv_{0}^{2}\left(1+\frac{4h^{2}}{l^{2}}\right)$。
}

\item
%% number: 22
%% paperName: 2003年上海高考第 22 题 14 分
%% typeId: 6
%% type: 计算题
%% chapter: 电磁感应
%% point: 电磁感应电路分析
%% method: 无
%% score: 14
%% degree: 450
%% duplicateId: 0
%% body:
如图所示，OACO 为置于水平面内的光滑闭合金属导轨，O、C 处分别接有短电阻丝（图中用粗线表示），$R_{1}=4\UO$、$R_{2}=8\UO$（导轨其它部分电阻不计），导轨 OAC 的形状满足方程 $y=2\sin\left(\frac{\pi}{3}x\right)$（单位：$\Um$），磁感强度 $B=0.2\UT$ 的匀强磁场方向垂直于导轨平面，一足够长的金属棒在水平外力 $F$ 作用下，以恒定的速率 $v=5\Ums$ 水平向右在导轨上从 O 点滑动到 C 点，棒与导轨接触良好且始终保持与 OC 导轨垂直，不计棒的电阻，求：

\begin{enumerate}
	\item
	外力 $F$ 的最大值；
	\item
	金属棒在导轨上运动时电阻丝 $R_{1}$ 上消耗的最大功率；
	\item
	在滑动过程中通过金属棒的电流 $I$ 与时间 $t$ 的关系。
\end{enumerate}

\onepicture[align=right, w=6.0cm]{\includesvg{figs/22}}

%% answer:
\jdanswer{
\begin{enumerate}
	\item
	$F_{\max}=0.3\UN$
	\item
	$P_{1}=1\UW$
	\item
	$I=\frac{3}{4}\sin\left(\frac{5\pi}{3}t\right)$
\end{enumerate}
}
%% memo:
\memoanswer{
（1）由题意知，金属棒做匀速运动，由平衡条件得 $F_{\text{外}}=F_{\text{安}}$。

金属棒切割磁感线产生的电动势为 $E=Blv$，回路中的电流为 $I=\frac{E}{R_{\text{总}}}$。

回路中的总电阻为 $R_{\text{总}}=\frac{R_{1}R_{2}}{R_{1}+R_{2}}=\frac{8}{3}\UO$。

金属棒受到的安培力大小为 $F_{\text{安}}=BIl=\frac{B^{2}l^{2}v}{R_{\text{总}}}$，故 $F_{\text{外}}=\frac{B^{2}l^{2}v}{R_{\text{总}}}$。

又 $l_{\max}=2\sin\frac{\pi}{2}=2\Um$，联立解得 $F_{\max}=0.2^{2}\times2^{2}\times5\times\frac{3}{8}\UN=0.3\UN$。

（2）由题意知，$l$ 取最大值时，电阻丝 $R_{1}$ 上消耗的功率最大，则有 $P_{1}=\frac{E^{2}}{R_{1}}=\frac{B^{2}l^{2}v^{2}}{R_{1}}=\frac{0.2^{2}\times2^{2}\times5^{2}}{4}\UW=1\UW$。

（3）金属棒与导轨接触点间的长度随时间的变化规律，有 $l=2\sin\frac{\pi}{3}x$，且 $E=Blv$，$x=vt$，解得 $I=\frac{E}{R_{\text{总}}}=\frac{Bv}{R_{\text{总}}}2\sin\left(\frac{\pi}{3}vt\right)=\frac{3}{4}\sin\left(\frac{5\pi}{3}t\right)$。
}

\item
%% number: 23
%% paperName: 2003年上海高考第 23 题 14 分
%% typeId: 6
%% type: 计算题
%% chapter: 电场
%% point: 带电粒子的直线运动
%% method: 极值
%% score: 14
%% degree: 250
%% duplicateId: 0
%% body:
为研究静电除尘，有人设计了一个盒状容器，容器侧面是绝缘的透明有机玻璃，它的上下底面是面积 $A=0.04\Umq$ 的金属板，间距 $L=0.05\Um$，当连接到 $U=2500\UV$ 的高压电源正负两极时，能在两金属板间产生一个匀强电场，如图所示，现把一定量均匀分布的烟尘颗粒密闭在容器内，每立方米有烟尘颗粒 $10^{13}$ 个，假设这些颗粒都处于静止状态，每个颗粒带电量为 $q=+1.0\times10^{-17}\UC$，质量为 $m=2.0\times10^{-15}\Ukg$，不考虑烟尘颗粒之间的相互作用和空气阻力，并忽略烟尘颗粒所受重力，求合上电键后：

\begin{enumerate}
	\item
	经过多长时间烟尘颗粒可以被全部吸收；
	\item
	除尘过程中电场对烟尘颗粒共做了多少功？
	\item
	经过多长时间容器中烟尘颗粒的总动能达到最大？
\end{enumerate}

\onepicture[align=right, w=6.0cm]{\includesvg{figs/23}}

%% answer:
\jdanswer{
\begin{enumerate}
	\item
	$t=0.02\Us$
	\item
	$W=2.5\times10^{-4}\UJ$
	\item
	$t_{1}=0.014\Us$
\end{enumerate}
}
%% memo:
\memoanswer{
（1）当最靠近上表面的烟尘颗粒被吸附到下板时，烟尘就全部吸附，烟尘颗粒所受电场力为 $F=\frac{qU}{L}$，由牛顿第二定律得 $a=\frac{F}{m}$，又 $L=\frac{1}{2}at^{2}=\frac{qUt^{2}}{2mL}$，联立解得 $t=L\sqrt{\frac{2m}{qU}}=0.02\Us$。

（2）由题意知，烟尘颗粒均匀分布在容器中，烟尘颗粒总数为 $NAL$，利用等效思想理解，全部颗粒等效为一个颗粒，且位于两金属板中间，带电量为 $NALq$，该颗粒运动到极板时，电场力做的功为 $W=\frac{1}{2}NALqU=2.5\times10^{-4}\UJ$。

（3）设某一烟尘颗粒下落的距离为 $x$，此时仍存在于容器中的颗粒数为 $NA(L-x)$，每个颗粒的动能为 $\frac{1}{2}mv^{2}$，由动能定理得 $\frac{qU}{L}x=\frac{1}{2}mv^{2}$。存在于容器中的颗粒数的总动能为 $E_{k}=\frac{1}{2}mv^{2}\cdot NA(L-x)=\frac{qU}{L}xNA(L-x)$。由函数知识知，当 $x=\frac{1}{2}L$ 时，$E_{k}$ 最大。由运动学公式得 $x=\frac{1}{2}at_{1}^{2}$，解得 $t_{1}=\sqrt{\frac{2x}{a}}=L\sqrt{\frac{m}{qU}}=0.014\Us$。
}

\end{enumerate}

\end{document}
